% BHU PYQ -- Manually transcribed & error-corrected
% Paper: GLB-604: Economic Geology
% Exam: B.Sc. (Hons.) Semester VI Examination, 2022-23

\documentclass[11pt,a4paper]{article}
\usepackage[a4paper,top=2.5cm,bottom=2.5cm,left=2.8cm,right=2.8cm,headheight=14pt]{geometry}
\usepackage{amsmath,amssymb}
\usepackage[T1]{fontenc}
\usepackage[utf8]{inputenc}
\usepackage{lmodern,microtype,fancyhdr,enumitem,hyperref,lastpage,parskip}

\hypersetup{colorlinks=true,urlcolor=black,linkcolor=black,
  pdftitle={GLB-604 Economic Geology -- BHU B.Sc. Sem VI 2022-23}}
\pagestyle{fancy}\fancyhf{}
\renewcommand{\headrulewidth}{0.4pt}\renewcommand{\footrulewidth}{0.4pt}
\fancyhead[L]{\small   \textit{Banaras Hindu University}}
\fancyhead[R]{\small   \textit{GLB-604: Economic Geology}}
\fancyfoot[C]{\small Page  \thepage\ of \pageref{LastPage}}


\newlist{parts}{enumerate}{2}
\setlist[parts,1]{label=(\alph*),leftmargin=*,itemsep=5pt,topsep=3pt}

\newcommand{\pts}[1]{\hfill   \textit{[#1]}}


\begin{document}

\begin{center}
  {\large\bfseries BANARAS HINDU UNIVERSITY}\\[2pt]
  {
\normalsize B.Sc. (Hons.) --- Semester VI Examination, 2022-23}\\[6pt]
  \rule{0.8\linewidth}{0.5pt}\\[6pt]
  {\large\bfseries Paper: GLB-604 --- Economic Geology}\\[4pt]
  {
\normalsize Time: 3 Hours \hfill Maximum Marks: 70}\\[2pt]
  {\small Ref. No.: 058443}
\end{center}
\rule{\linewidth}{0.4pt}
\smallskip



\noindent   \textit{Instructions: Answer   \textbf{five} questions in all, selecting at least   \textbf{two} each from Sections A and B, and Section-C which is compulsory. The figures in the right-hand margin indicate marks.}
\bigskip

\bigskip


\noindent {\large\bfseries SECTION --- A}
\rule{\linewidth}{0.4pt}\smallskip




\noindent   \textbf{Question 1.} What is cavity filling? Explain the characteristic features of cavity filling with suitable sketches. \pts{14}

\medskip


\noindent   \textbf{Question 2.} Write about the origin, modes of occurrence, and Indian distribution of any two minerals used in the chemical industry. \pts{14}

\medskip


\noindent   \textbf{Question 3.} Write short notes on   \textbf{any two} of the following: \pts{$2  \times 7 = 14$}

\begin{parts}
  \item Classification of gemstones
  \item Raw materials for phosphate in the fertilizer industry
  \item Residual liquid injection
\end{parts}

\medskip


\noindent   \textbf{Question 4.} What is a placer deposit? What are placer minerals? Discuss the formation of placer deposits. \pts{14}

\bigskip


\noindent {\large\bfseries SECTION --- B}
\rule{\linewidth}{0.4pt}\smallskip




\noindent   \textbf{Question 5.} Define ore body in a mineral deposit. Describe various morphologies of ore bodies in mineral deposits with suitable diagrams. \pts{14}

\medskip


\noindent   \textbf{Question 6.} Describe the mineralogy, modes of occurrence, economic uses, and distribution of iron ore deposits in India. \pts{14}

\medskip


\noindent   \textbf{Question 7.} What is petroleum? Briefly describe the geology, modes of occurrence, and distribution of petroleum deposits in India. \pts{14}

\bigskip


\noindent {\large\bfseries SECTION --- C}
\rule{\linewidth}{0.4pt}\smallskip




\noindent   \textbf{Question 8.} Answer the following:

\begin{parts}
  \item What is paragenesis? \pts{2}
  \item Define tenor of ore with an example. \pts{2}
  \item Define the physical constituents of coal. \pts{2}
  \item Write the name of an important ore mineral of tin. \pts{1}
  \item What is gossan? \pts{2}
  \item Write the names of any two organic materials used as gems. \pts{2}
  \item What is refractory? \pts{2}
  \item Why is wollastonite used in the ceramic industry? \pts{1}
\end{parts}

\vfill\rule{\linewidth}{0.4pt}

\begin{center}\small
\href{https://bhu-exam.vercel.app/}{Click here to visit our website}\\[4pt]
\href{https://me-aryan.vercel.app/}{Click here to visit our contributor}\\[4pt]
\href{https://quashan-me.vercel.app/}{Click here to visit our contributor}
\end{center}
\end{document}
